\section{Discussion}

\subsection{Interpreting Sparse Coupling Detection Capability}

Our core finding---methodology can detect coupling limited to 2 dominant dimension pairs rather than distributed across all 10 pairs---demonstrates the framework's ability to distinguish sparse from broad patterns. If validated on real benchmarks, sparse coupling would suggest trustworthiness dimensions are \textbf{architecturally independent by default}. Coupling would emerge only where specific mechanisms overlap: calibration failures (truthfulness-robustness) and value alignment training (fairness-safety).

This capability would help refute two alternative hypotheses. First, the ``broad independence'' hypothesis predicts zero coupling across all dimension pairs, assuming trustworthiness dimensions reflect orthogonal model capabilities. Our synthetic validation shows coupling is detectable when present (6 pairs, phi 0.33--0.40, $p < 1\times10^{-13}$). Second, the ``broad coupling'' hypothesis predicts pervasive co-occurrence across many pairs, assuming shared architectural bottlenecks cause failures to cascade across dimensions. Our h-m2 FAIL result demonstrates methodology can detect sparsity: 0 models exhibit $\geq$3 significant pairs in synthetic validation.

\textbf{Why sparse coupling detection matters if validated:} It would reveal that multi-dimensional trustworthiness failures are not universal---models do not simply ``fail everywhere when stressed.'' Instead, vulnerabilities would cluster in specific subsystems detectable by this framework. Truthfulness and robustness would couple because both depend on calibration and factual grounding; when adversarial prompts exploit calibration weaknesses, both dimensions fail simultaneously. Fairness and safety would couple because both are shaped by value alignment training (RLHF); correlated reward model biases drive co-occurrence in these dimensions. The remaining 8 dimension pairs would remain largely independent, indicating separate processing pathways.

\subsection{Difficulty as Suppressor in Synthetic Data---Alternative Explanations Untested}

The h-m1 finding---effect size retention 86--161\%, with 5 of 6 cases showing partial phi exceeding raw phi---demonstrates methodology capability to control for difficulty confounds. However, the suppressor effect interpretation (difficulty masks true coupling) may not replicate with real benchmarks. Prior work predicts confounds should inflate raw effect sizes; controlling for confounds should reduce estimates by 40--60\%. Our synthetic results show the opposite: controlling difficulty \textbf{strengthens} coupling estimates.

We interpret difficulty as a \textbf{suppressor variable} in synthetic data: instances failing due to high difficulty may pass on dimension-specific vulnerabilities, and vice versa. Difficulty variance adds noise orthogonal to coupling signal; removing this noise unmasks latent coupling strength. This aligns with our independence constraint ($|\text{corr}(\text{difficulty}, \text{dimension})| < 0.2$)---difficulty was forced to be uncorrelated with trustworthiness dimensions by design.

\textbf{Alternative explanations remain untested:}
\begin{enumerate}
\item \textbf{Synthetic constraint artifact:} Difficulty scores were constrained $|r| < 0.2$ by design to meet partial correlation independence assumptions. Real benchmark difficulty distributions may show different relationships.
\item \textbf{Statistical artifact:} Partial correlation may inflate effect sizes in small samples with weak predictors, independent of true suppressor mechanisms.
\item \textbf{Measurement error:} Confidence-based difficulty proxy (simulated in Phase 4) may not capture true difficulty as measured by real benchmark metadata.
\end{enumerate}

Real-world validation with unconstrained difficulty distributions is required to confirm this mechanism. If suppressor effect replicates, it would have methodological implications for multi-dimensional LLM evaluation---difficulty control may \textbf{reveal} stronger latent relationships rather than reduce artifacts. If it doesn't replicate, partial correlation remains useful as confound control (preventing inflation), even without suppressor effects.

\subsection{Honest Limitations}

\textbf{L1: Synthetic data (HIGH IMPACT).} All Phase 4 experiments used synthetic coupling data instead of real benchmarks (MultiTrust/TrustLLM remain gated). We validated the measurement methodology---phi coefficient, partial correlation, Mantel test implementations work correctly---but cannot claim these coupling patterns exist in real LLMs. Synthetic data was designed with target coupling values (phi 0.35--0.40 for dominant pairs), so results confirm ``coupling is measurable if it exists'' rather than ``coupling exists in practice.''

\textbf{Why this is acceptable for a proof-of-concept:} The contribution is methodological---we demonstrate a pipeline for detecting and characterizing coupling. Real-world characterization is the natural next step (Phase 5 baseline comparison), but the technique's feasibility is established. Analogously, early benchmark papers (GLUE, SuperGLUE) validated evaluation frameworks on preliminary datasets before scaling to production models.

\textbf{Future mitigation:} Phase 5 MUST use real MultiTrust/TrustLLM data with GPT-4, Claude-3, Llama-3 API evaluations to upgrade findings from ``methodology demonstration'' to ``empirical characterization.''

\textbf{L2: Sample size (MEDIUM IMPACT).} Mantel test for h-c1 model-specific fingerprints used n=100 instances/dimension, insufficient for statistical power (requires $n \geq 500$ per Legendre \& Legendre power analysis). Criterion $r < 0.7$ was met (coupling matrices structurally dissimilar), but p-values remained non-significant (0.317--0.758 $\gg$ 0.0167).

\textbf{Why this is acceptable:} Qualitative patterns in synthetic data are consistent---distinct coupling patterns observable across GPT-4 (cognitive coherence chain), Claude-3 (value alignment cluster), and Llama-3 (minimal coupling). The statistical inconclusiveness reflects sample size limitation, not absence of effect. This is a standard Phase 4 PoC trade-off: demonstrate feasibility at small scale, power appropriately for confirmatory Phase 5.

\textbf{Future mitigation:} Increase to $n \geq 1000$ instances/model (200/dimension) or leverage real benchmarks with broader coverage (MultiTrust 8 dimensions $\times$ 1000 instances).

\textbf{L3: Bonferroni over-correction (LOW IMPACT).} h-m2 applied strict Bonferroni correction (alpha\_adj = 0.01/30 $\approx$ 0.00033) to control family-wise error rate across 30 tests. This excludes borderline pairs (e.g., GPT-4 truthfulness-robustness phi 0.302, p\_adj=0.057) that might achieve significance under less conservative methods (FDR control).

\textbf{Why this is acceptable:} Sparse coupling conclusion is robust to correction method choice. Even if 1--2 additional pairs achieve significance under FDR, no model would reach $\geq$3 pairs threshold. The literature independently supports sparse coupling (Li \& Li 2024 triangular trade-offs document specific robustness-fairness couplings, not broad patterns). Bonferroni conservatism protects against false positives; the trade-off (potentially missing weak couplings) is acceptable for establishing core sparsity capability.

\textbf{Future sensitivity analysis:} Report both Bonferroni and FDR-corrected results to demonstrate robustness; explore phi 0.20--0.29 range for additional weak couplings.

\subsection{Broader Impact}

\textbf{Positive applications if validated:} Coupling-aware model selection could reduce compound failures in safety-critical deployments. Deployment teams could prioritize models with strong coupling in required dimensions---e.g., select Claude for fairness-safety critical scenarios (based on synthetic profile patterns) or GPT-4 for truthfulness-robustness requirements. Benchmark designers could ensure coverage of detected coupling pairs (truth-robust, fair-safe) for compound risk assessment rather than treating dimensions independently.

\textbf{Potential misuse:} Adversarial actors might exploit coupling patterns to maximize attack impact---targeting truthfulness knowing robustness will fail simultaneously. However, coupling as a transparency tool (publishing coupling profiles alongside model cards) is preferable to security-through-obscurity. Defenders can use coupling knowledge to design targeted mitigations (e.g., calibration-focused interventions to address truth-robust cluster).

\textbf{Limitations as guardrails:} The synthetic data limitation (L1) prevents premature deployment recommendations until real-world coupling patterns are validated. This work serves as methodology demonstration rather than actionable guidance. This is appropriate for a foundational study---establish measurement approach before prescribing interventions.

\subsection{What We Learned}

Sparse coupling detection (methodology distinguishes 2 pairs from 10) is a feasible measurement capability, validated on synthetic data with difficulty-independence controls. If replicated on real benchmarks, it would refine the multi-dimensional evaluation paradigm: dimensions are neither universally independent (some coupling detectable) nor universally coupled (most pairs remain uncorrelated). The specific vulnerability clusters---calibration failures (truth-robust) and value alignment (fair-safe)---would suggest targeted rather than universal mitigation strategies.

The difficulty suppressor effect in synthetic data demonstrates control methodology but requires real-world validation with unconstrained difficulty distributions. Researchers should report both raw and difficulty-controlled effect sizes, as control may reveal rather than reduce signal strength (if suppressor effect replicates) or prevent inflation (if it doesn't).

Model-specific fingerprints remain observational but statistically unconfirmed pending larger samples. The patterns in synthetic data (GPT-4 cognitive coherence, Claude-3 value alignment, Llama-3 independence) are consistent across metrics and would align with known architectural/training differences if validated, warranting follow-up with appropriately powered samples.
